\documentclass{article}
\usepackage{hyperref}
\usepackage{amsmath}
\begin{document}
This readme is meant to explain where the system of differential equations implimented in the julia script came from. If the simulation is faulty, the bug may be in the derivation of the system, so it's necessary to preserve it.

The system simulates a multilayer railgun (my own improvement on an augmented railgun). A railgun moves a projectile by passing current from one rail through the projectile into the other rail. An augmented railgun increases (or "augments") the magnetic field by an electromagnet. A multilayer railgun uses each winding in the electro magnet as it's own rail, and passes current through a new layer of the projectile.

The power comes from a capacitor of voltage $V_s$. The magnetic field is generated by the rails in combination with permanent magnets. Other components, such as the thyristor, are not included in the simulation, except that they contribute to $R$, the sum of all resistances. The current, $I$, passes through the projectile between the rails, and is assumed to be the same everywhere in the circuit.

The total magnetic flux in the circuit is denoted as $\phi$. The changing current in the rails and the projectile moving through magnetic field are both a source of change in flux, $\frac{d\phi}{dt}$.

The following equation can be found \href{https://physics.stackexchange.com/questions/791718/how-to-calculate-impedance-in-a-wire-moving-through-a-magnetic-field}{here}.

\[V_s-V_b=IR\]

However, the above link says $V_b$ depends on the velocity of the projectile, B-field, and projectile length in the $y$ direction. In reality, $V_b$ changes with any change in magnetic flux in the circuit. For this line of reasoning we will use two sources for change in flux (from the changing magnetic field, and from the changing size of the circuit). Read more \href{https://www.maxwells-equations.com/faraday/faradays-law.php}{here}

\[V_b=\frac{d\phi}{dt}\]
\[\phi=N\int\int_0^d Bdxdy\]

The $x$ direction is parallel to the rails, $X$ is the length of the rails and $d$ is the distance of the projectile down the rails.

\[\phi=YN\int_0^d Bdx\]

The $y$ direction is perpendicular to the rails and the magnetic field and $Y$ is the length of the rails in the $y$ direction. We assume the B field is uniform in the Y direction, for the sake of simplicity.

the $H$ field (and by extension the $B$ field) is generated by current in the rails, and large permanent magnets. \href{https://pressbooks.online.ucf.edu/osuniversityphysics2/chapter/magnetic-field-due-to-a-thin-straight-wire/}{This} link gives the following equation, where $r$ is the distance between the wire and the center of the rails in the $y$ direction, and $l$ is the signed distance in the $x$ direction between the point of interest and the current carrying segment of wire.

\[B(x) = \frac{\mu I}{4\pi} \int_{wire}\frac{r}{(l^2+r^2)^{3/2}}dl\]

This equation is for one rail. We will double it for one full layer. Since the current \emph{mostly} run in the rails, we will simplify the equations by only integrating across the $x$ direction (hence doubling: once there, once back, no side to side). We also will give it a subscript $n$, because it is the $n$th of $N$ rails. Each rail will be some distance $r$ from the center of the rails, so $r$ is a function of $n$.

\[B = \frac{\mu I}{2\pi} \int_{-x}^{d-x}\frac{r}{(l^2+r^2)^{3/2}}dl\]

The total magnetic field is the sum of each layer, plus the magnetic field due to the permanent magnets $\mu H$:

\[B = \sum_{n=0}^{N}B_n + \mu H = \frac{\mu I}{2\pi} \left(\sum_{n=1}^{N} \int_{x=0}^{x=d}\frac{r}{(l^2+r^2)^{3/2}}dl\right) + \mu H\]

Nothing inside the sumation varies with time, so it makes conceptual sense to seperate it out, and it avoids repetition later on. We will call this value... $M$ for "b field modifier"????

\[M_n=\int_{x=0}^{x=d}\frac{r}{(l^2+r^2)^{3/2}}dl = \left .\frac{l}{r\sqrt{r^2+l^2}} \right\rvert_{x=0}^{x=d}\]

$l$ in this equation is the distance from the point of interest, rather than the origin. (a negative value of $l$ is behind the point of interest). But the origin of $x$ is the beginning of the rails. If the the point of interest is at some value of $x$, then the origin of $x$ axis is $-x$ in terms of $l$.

\[M_n=\left .\frac{l}{r\sqrt{r^2+l^2}} \right\rvert_{l=-x}^{l=d-x} = r^{-1}\left (\frac{d-x}{\sqrt{r^2+(d-x)^2}} + \frac{x}{\sqrt{r^2+x^2}}\right )\]
\[M = \sum_{n=0}^{N} M_n \]
\[B=\frac{\mu I}{2\pi}M + \mu H\]

\[ \int Bdx = \frac{\mu I}{2\pi} \int \sum_0^N M_ndx +\mu Hx\]
\[ = \frac{\mu I}{2\pi} \int \sum r^{-1}\left (\frac{d-x}{\sqrt{r^2+(d-x)^2}} + \frac{x}{\sqrt{r^2+x^2}}\right )dx +\mu Hx\]
\[ = \frac{\mu I}{2\pi} \sum r^{-1} ( \sqrt{r^2 + x^2} - \sqrt{r^2 + (d-x)^2}) + \mu Hx + K \]

Now that we have the indefinite integral, we will evaluate it on it's bound, from $0$ to $d$.

\[ \int_0^d Bdx = \frac{\mu I}{2\pi} \sum \left[r^{-1} ( \sqrt{r^2 + d^2} - \sqrt{r^2 + (d-d)^2}) - r^{-1} ( \sqrt{r^2 + 0^2} - \sqrt{r^2 + (d-0)^2})\right] + \mu Hd \]
\[ = \frac{\mu I}{2\pi} \left(\sum r^{-1}(2\sqrt{r^2 + d^2} - 2r) \right) + \mu Hd\]

We will plug these values in for the equation for $\phi$, which, as a reminder is: 

\[\phi=YN\int_0^d Bdx\]
\[ = YN\frac{\mu I}{2\pi } \left(\sum r^{-1}(2\sqrt{r^2 + d^2} - 2r) \right) + \mu Hd\]

Remember, the reason we're interested in $\phi$ is because it's derivative appears in this equation.

\[V_b=\frac{d\phi}{dt}\]

So we take the derivative. By the multiplication rule we have:
\[ \frac{d\phi}{dt} = I^\prime \frac{YN\mu}{2\pi} \left(\sum r^{-1} (2\sqrt{r^2 + d^2} - 2r) \right) + I \frac{YN\mu}{2\pi} \frac{d}{dt}\left(\sum r^{-1} (2\sqrt{r^2 + d^2} - 2r) \right) + \mu H Y d^\prime \]

\[ = I^\prime \frac{YN\mu}{2\pi} \left(\sum r^{-1} (2\sqrt{r^2 + d^2} - 2r) \right) + I \frac{YN\mu}{2\pi} \left(\sum r^{-1} (\frac{2dd^\prime}{\sqrt{r^2 + d^2}}) \right) + \mu H Y d^\prime \]

Now to combine like terms and solve for $I^\prime$.

\[I^\prime = \frac{  \frac{d\phi}{dt} - I \frac{YN\mu}{2\pi} \sum r^{-1} (\frac{2dd^\prime}{\sqrt{r^2 + d^2}}) - \mu H Y d^\prime }{ \frac{YN\mu}{2\pi} \sum r^{-1} (2\sqrt{r^2 + d^2} - 2r) } \]

And finally remember that $\frac{d\phi}{dt} = V_b$, and $V_s - V_b = IR$  or $\frac{d\phi}{dt} = V_s - IR $

Now that we have an equation for $I^\prime$ we can talk about $V_s$, the source voltage. The source voltage is due to a capacitor. Capacitors decay exponentially to an input voltage. I.e. if we have the capacitor at 100 volts, and we're "charging" it at 80 volts, it will decay to 80 volts. You can read more \href{https://physics.stackexchange.com/questions/800396/discharge-of-a-capacitor-with-variable-backvoltage}{here}, where I got the following equation.

\[RC\frac{dV_s}{dt} + V_s = V_b \]
\[\frac{dV_s}{dt}= \frac{V_b - V_s}{RC} = \frac{-I}{C}\]

And we have one last variable to talk about, $\frac{d^2d}{dt^2}$. How the heck is the projectile going to start moving? from wikipedia article on the lorentz force we have:

\[ f = Il\times B \]

This entire time we've been thinking about things as scalars, and I sure as heck ain't gonna change that now. So we will assume the $B$ vector and $Il$ vectors are perfectly perpendicular. It should be noted that the $l$ in the equation from wikipedia is what we've been calling $w$, the width of the projectile, times the number of layers $N$. $B$ is of course a function of $x$, and in this case, it's when $x=d$, the distance of the projectile down the rails.

\[f = INwB(d) \]
\[f = I^2Nw\frac{\mu}{2\pi}M(d) + INw\mu H \]
\[f = I^2Nw\frac{\mu}{2\pi}\left(\sum_0^N \frac{d}{r\sqrt{r^2-d^2}}\right) + INw\mu H \]

Here $f$ is the lorentz force. The total forces acting on the projectile are $f_{net}=f-f_s$ where $f_s$ is the force due to friction against the rails.

\[f_{net}=m\frac{dd^2}{d^2t} \]
\[\frac{dd^2}{d^2t} = m^{-1}\left(I^2w\frac{\mu_p}{2\pi}\left(\sum_0^N \frac{d}{r\sqrt{r^2-d^2}}\right) + Iw\mu H - f_s\right) \]

Finally, we need an equation for $r$, which you'll recall is a function $r(n)$, the distance between the $n$th rail and the center of the bore. While this deserves a more thorough treatment, for now I'm going to assume it is a constant distance.

And... I think that's it. That's everything. There should now be an equation for each of the four output variables now (Remember, the equation for $d'$ is $d'=d'$. If you're wondering why that's valid here, but I can't just say $v_s = v_s$, it's because $d'$ is already in the input vector). It's time to implement this in the code! Good luck dear reader!

\end{document}